\documentclass[10pt,twocolumn]{article}
\usepackage[margin=0.8in]{geometry}
\usepackage[T1]{fontenc}
\usepackage{lmodern}
\usepackage{booktabs}
\usepackage{amsmath}
\usepackage{microtype}
\usepackage[hidelinks]{hyperref}
% Generated by make_numbers.py from results_v2.json. Do not edit.
\newcommand{\ablChalClfF}{90.2}
\newcommand{\ablChalClfFpr}{0.0}
\newcommand{\ablChalFullF}{96.0}
\newcommand{\ablChalFullFpr}{0.0}
\newcommand{\ablChalIntentsF}{96.0}
\newcommand{\ablChalIntentsFpr}{0.0}
\newcommand{\ablChalProtF}{96.0}
\newcommand{\ablChalProtFpr}{0.0}
\newcommand{\ablChalRulesF}{85.2}
\newcommand{\ablChalRulesFpr}{29.1}
\newcommand{\ablHinBClfF}{91.3}
\newcommand{\ablHinBClfFpr}{0.0}
\newcommand{\ablHinBFullF}{100.0}
\newcommand{\ablHinBFullFpr}{0.0}
\newcommand{\ablHinBIntentsF}{100.0}
\newcommand{\ablHinBIntentsFpr}{0.0}
\newcommand{\ablHinBProtF}{100.0}
\newcommand{\ablHinBProtFpr}{0.0}
\newcommand{\ablHinBRulesF}{67.6}
\newcommand{\ablHinBRulesFpr}{96.0}
\newcommand{\ablRealClfF}{98.1}
\newcommand{\ablRealClfFpr}{2.7}
\newcommand{\ablRealFullF}{97.5}
\newcommand{\ablRealFullFpr}{3.9}
\newcommand{\ablRealIntentsF}{97.7}
\newcommand{\ablRealIntentsFpr}{4.1}
\newcommand{\ablRealProtF}{97.7}
\newcommand{\ablRealProtFpr}{3.9}
\newcommand{\ablRealRulesF}{79.1}
\newcommand{\ablRealRulesFpr}{90.8}
\newcommand{\accRealDep}{96.3}
\newcommand{\accRealRep}{97.1}
\newcommand{\accRealTwo}{96.7}
\newcommand{\aucRealTwo}{0.994}
\newcommand{\bandAccOut}{97.7}
\newcommand{\bandErrIn}{20}
\newcommand{\bandErrTotal}{61}
\newcommand{\bandShare}{3.6}
\newcommand{\chgCombo}{78.4}
\newcommand{\chgDefang}{43.2}
\newcommand{\chgHomo}{76.2}
\newcommand{\chgLeet}{77.1}
\newcommand{\chgSpaced}{78.0}
\newcommand{\chgZero}{78.0}
\newcommand{\ciHiRHDep}{75.6}
\newcommand{\ciHiRHRep}{80.6}
\newcommand{\ciHiRHTwo}{78.3}
\newcommand{\ciHiRealDep}{97.8}
\newcommand{\ciHiRealRep}{98.3}
\newcommand{\ciHiRealTwo}{98.0}
\newcommand{\ciLoRHDep}{50.0}
\newcommand{\ciLoRHRep}{55.6}
\newcommand{\ciLoRHTwo}{52.3}
\newcommand{\ciLoRealDep}{96.4}
\newcommand{\ciLoRealRep}{97.1}
\newcommand{\ciLoRealTwo}{96.8}
\newcommand{\eceClf}{0.008}
\newcommand{\eceFinal}{0.017}
\newcommand{\fChalDep}{96.9}
\newcommand{\fChalRep}{95.6}
\newcommand{\fChalTwo}{96.0}
\newcommand{\fHinADep}{88.8}
\newcommand{\fHinARep}{85.7}
\newcommand{\fHinATwo}{95.5}
\newcommand{\fHinBDep}{85.0}
\newcommand{\fHinBRep}{85.2}
\newcommand{\fHinBTwo}{100.0}
\newcommand{\fRHDep}{63.9}
\newcommand{\fRHRep}{69.6}
\newcommand{\fRHTwo}{66.7}
\newcommand{\fRealDep}{97.1}
\newcommand{\fRealRep}{97.7}
\newcommand{\fRealTwo}{97.5}
\newcommand{\fprAdvCleanOne}{3.9}
\newcommand{\fprAdvCleanTwo}{3.9}
\newcommand{\fprAdvComboOne}{2.4}
\newcommand{\fprAdvComboTwo}{3.9}
\newcommand{\fprAdvDefangOne}{3.2}
\newcommand{\fprAdvDefangTwo}{3.9}
\newcommand{\fprAdvHomoOne}{3.0}
\newcommand{\fprAdvHomoTwo}{3.9}
\newcommand{\fprAdvLeetOne}{3.2}
\newcommand{\fprAdvLeetTwo}{3.9}
\newcommand{\fprAdvSpacedOne}{3.5}
\newcommand{\fprAdvSpacedTwo}{3.8}
\newcommand{\fprAdvZeroOne}{2.7}
\newcommand{\fprAdvZeroTwo}{3.9}
\newcommand{\fprChalDep}{0.0}
\newcommand{\fprChalRep}{0.0}
\newcommand{\fprChalTwo}{0.0}
\newcommand{\fprHinADep}{9.3}
\newcommand{\fprHinARep}{9.3}
\newcommand{\fprHinATwo}{9.3}
\newcommand{\fprHinBDep}{12.7}
\newcommand{\fprHinBRep}{6.0}
\newcommand{\fprHinBTwo}{0.0}
\newcommand{\fprRHDep}{33.3}
\newcommand{\fprRHRep}{23.8}
\newcommand{\fprRHTwo}{26.2}
\newcommand{\fprRealDep}{3.9}
\newcommand{\fprRealRep}{3.2}
\newcommand{\fprRealTwo}{3.9}
\newcommand{\grpRHDep}{60.0}
\newcommand{\grpRHRep}{66.7}
\newcommand{\grpRHTwo}{64.2}
\newcommand{\linkCrossF}{94.1}
\newcommand{\linkPhoneRealF}{69.3}
\newcommand{\linkRealF}{56.2}
\newcommand{\linkRealFpr}{36.1}
\newcommand{\msPerMsg}{0.71}
\newcommand{\nChal}{259}
\newcommand{\nDatedCamps}{2551}
\newcommand{\nFutureCamps}{221}
\newcommand{\nFutureMsgs}{855}
\newcommand{\nHinTrain}{1{,}200}
\newcommand{\nNoLink}{642}
\newcommand{\nRH}{77}
\newcommand{\nRHGen}{42}
\newcommand{\nRHGroups}{60}
\newcommand{\nRHScam}{35}
\newcommand{\nReal}{1874}
\newcommand{\nTrain}{9{,}513}
\newcommand{\pHinB}{$<$0.001}
\newcommand{\pRH}{0.38}
\newcommand{\pReal}{0.12}
\newcommand{\precRHDep}{62.2}
\newcommand{\precRHRep}{70.6}
\newcommand{\precRHTwo}{67.6}
\newcommand{\recCleanOne}{96.4}
\newcommand{\recCleanTwo}{97.1}
\newcommand{\recComboOne}{88.4}
\newcommand{\recComboTwo}{97.2}
\newcommand{\recDefangOne}{96.0}
\newcommand{\recDefangTwo}{97.1}
\newcommand{\recFuture}{99.6}
\newcommand{\recHomoOne}{94.4}
\newcommand{\recHomoTwo}{97.1}
\newcommand{\recLeetOne}{94.5}
\newcommand{\recLeetTwo}{97.2}
\newcommand{\recNoLink}{95.8}
\newcommand{\recRHDep}{65.7}
\newcommand{\recRHRep}{68.6}
\newcommand{\recRHTwo}{65.7}
\newcommand{\recRandMax}{99.0}
\newcommand{\recRandMean}{97.8}
\newcommand{\recRandMin}{95.4}
\newcommand{\recSpacedOne}{94.3}
\newcommand{\recSpacedTwo}{96.9}
\newcommand{\recZeroOne}{94.0}
\newcommand{\recZeroTwo}{97.1}
\newcommand{\vocabTwo}{36{,}641}


\title{Hardening an Indian Scam-Message Detector:\\Evasion Robustness, a Train--Serve Mismatch,\\and a Synthetic-to-Real Gap in Hindi}
\author{Priyam Thakur \and Bhavya Goyal \and Samir Chaudhary \and Sana Subodh \and Kayarvizhy N\\[2pt]
\small Department of Computer Science and Engineering, BMS College of Engineering, Bengaluru, India}
\date{}

\begin{document}
\maketitle

\begin{abstract}
ScamGuard is a deployed detector for scam SMS and chat messages in India that combines a TF-IDF logistic-regression
classifier with interpretable intent rules. This paper reports what changed, and what did not, when we hardened it.
First, scammers disguise trigger words with lookalike letters, zero-width characters, spaced-out letters, leetspeak
and defanged links. Under these attacks, recall on \nReal{} held-out real messages falls from \recCleanOne\% to
\recComboOne\%. A deterministic de-obfuscation step restores it to \recComboTwo\%, and such tricks already occur in the
real data. Second, an audit of the serving path found a train--serve mismatch: the live app appended an English concept
summary that the model never saw in training, which lowered F1 from \fRealRep\% to \fRealDep\% and raised false alarms
from \fprRealRep\% to \fprRealDep\%. Third, after adding Hindi and Hinglish data and rules, the system scored
\fHinBTwo\% F1 on a held-out synthetic Hinglish set. On \nRH{} real Hindi messages labelled before scoring, it
scored only \fRHTwo\% F1 with \fprRHTwo\% false alarms, not significantly better than before ($p=\pRH$). We also
report a temporal test, calibration (ECE \eceFinal), an ``uncertain'' band that holds \bandShare\% of messages and
\bandErrIn{} of \bandErrTotal{} errors, a fusion ablation, and simple-rule baselines. Every number is produced by one
script that runs in under a minute on a laptop CPU.
\end{abstract}

\section{Introduction}
Fraud by SMS and messaging apps is common in India, and detectors trained on public English corpora transfer poorly
to Indian messages. ScamGuard~\cite{scamguard1} addressed this with ScamGuard-IN-2026, a corpus of \nTrain{}
training and development messages and \nReal{} real test messages from campaigns held out of training. It also adds
intent rules that explain each verdict. This paper studies three questions that a deployed detector faces after its
first evaluation:
\begin{enumerate}\itemsep0pt
\item Does it survive the cheap disguises that scammers use?
\item Does the live system behave like the evaluated one?
\item Do gains on synthetic code-mixed data carry over to real Hindi messages?
\end{enumerate}

\paragraph{Contributions.}
\begin{itemize}\itemsep0pt
\item An evasion benchmark of five character-level attacks and their combination, applied to the real test set.
  A rule-based de-obfuscator raises recall under the combined attack from \recComboOne\% to \recComboTwo\% at
  unchanged false alarms on clean text.
\item A documented train--serve mismatch in the deployed system, its measured cost, and its fix.
\item A cautionary result: \fHinBTwo\% F1 on a template-disjoint synthetic Hinglish test, but \fRHTwo\% on real Hindi
  messages labelled before scoring. We give an error analysis and release the real test set.
\item Supporting analyses (temporal split, calibration, review band, fusion ablation, rule baselines) and a one-command
  ingest path for newly collected real messages.
\end{itemize}

\section{System and data}
\paragraph{Classifier and risk engine.} The classifier is logistic regression ($C=16$, balanced weights) over word
1--2-gram and character 2--5-gram TF-IDF (\vocabTwo{} features in v2). A risk engine adds weighted evidence from
intent patterns (for example ``asks for OTP'', ``pay a fee to receive'', ``legal threat''). It subtracts protective
signals (``never share your OTP'', official helplines) in log-odds space and caps clearly genuine transaction
notices~\cite{scamguard1}.

\paragraph{Splits.} ScamGuard-IN-2026 groups near-duplicate messages into campaigns with MinHash. Whole campaigns are
assigned to train, development or test, so no campaign crosses the boundary. The real test set holds \nReal{}
messages; a hand-written challenge set holds \nChal{} hard cases.

\paragraph{Synthetic Hindi/Hinglish banks.} We wrote three template banks separately: TRAIN, TEST~A and TEST~B. TRAIN
is split by template, 80\% for training (\nHinTrain{} rows) and 20\% for development. Settings were chosen only on the
development split and the development templates. Set A was inspected once for error analysis, so we report it only
as a development check. Set B was written afterwards and scored once. No message text is shared between training and
either test bank.

\paragraph{Real Hindi test set.} IMC'25~\cite{imc25} contains user-reported smishing in many languages. We took the
127 Hindi reports that were not in ScamGuard-IN-2026. Each was labelled with the dataset's rule before any model
scored it. A message is a scam if it asks for a secret, a payment or fee, an app install, or a call or reply to a
private number, or if it pushes a non-official link with a lure. Fifty chat fragments, political or entertainment
promotions and unclear reports were dropped rather than guessed. The remaining \nRH{} messages (\nRHScam{} scam,
\nRHGen{} genuine) form \nRHGroups{} near-duplicate groups (character 5-gram Jaccard $\ge 0.6$). We report
message-level and group-level scores. One annotator labelled the set, so it has no agreement statistic.

\paragraph{Dates.} IMC'25 keeps the timestamp visible on each reported screenshot. Taking only explicit four-digit
years dates the first appearance of \nDatedCamps{} campaigns.

\section{Methods}
\paragraph{De-obfuscation.} A deterministic function runs before both the classifier and the rules. It applies NFKC
normalisation, removes zero-width characters, maps Cyrillic and Greek lookalikes to Latin, rewrites defanged links
(``hxxp'', ``[.]'', ``dot com''), and joins runs of single spaced or dotted letters. It also reverses leetspeak
inside words that are at least half letters, while keeping units, codes and ordinals (``4G'', ``10am'', ``2nd'',
``XX1234''). It never translates, and it leaves Devanagari untouched.

\paragraph{Hindi and Hinglish rules.} We added Devanagari and Hinglish versions of existing intents. Contact with a
private number counts only together with another intent, and a negation guard stops ``OTP kisi ko na batayein''
(``do not tell anyone the OTP'') from reading as a request.

\paragraph{Serving path.} The v1 app appended an English list of detected Hindi concepts to the message before scoring.
The model was trained on raw text, so this created a mismatch between training and serving~\cite{sculley2015}. v2
keeps the list for display only. The model input is the de-obfuscated text, exactly as in training.

\paragraph{Systems compared.} \emph{v1 reported} is the exam-version model and risk engine, scored on raw text as in
the v1 report. \emph{v1 deployed} is the same model scored the way the live v1 app scored it. \emph{v2} is trained on
the same data plus the synthetic Hinglish rows, with de-obfuscation and the new rules.

\paragraph{Attacks.} Each attack rewrites about 30 common trigger words (``kyc'', ``otp'', ``account'',
``verify'', \dots) in every message:
\begin{itemize}\itemsep0pt
\item \emph{homoglyph}: Cyrillic letters, each with probability 0.7;
\item \emph{zero-width}: U+200B between letters;
\item \emph{spaced}: ``c l i c k'';
\item \emph{leetspeak}: o$\to$0, e$\to$3, a$\to$4, i$\to$1, s$\to$5;
\item \emph{defanged URL}: ``hxxp'' and ``[.]'' (applied to links rather than words);
\item \emph{combined}: a random word-level trick plus defanging.
\end{itemize}
Attacks are applied to both classes, so false alarms are measured too.

\paragraph{Statistics.} We use 1{,}000-sample bootstrap 95\% confidence intervals, exact McNemar tests on paired
verdicts~\cite{mcnemar1947} and expected calibration error over 15 bins~\cite{guo2017}.

\section{Results}
\begin{table}[t]\centering\small
\setlength{\tabcolsep}{4pt}\caption{Final verdicts (classifier + risk engine). F1 and false-alarm rate (FPR) in \%. McNemar compares v1 deployed with v2.}
\label{tab:main}
\begin{tabular}{@{}lcccc@{}}\toprule
Test set & v1 rep. & v1 dep. & v2 & $p$ \\\midrule
Real test, F1 & \fRealRep & \fRealDep & \fRealTwo & \pReal \\
Real test, FPR & \fprRealRep & \fprRealDep & \fprRealTwo & \\
Challenge, F1 & \fChalRep & \fChalDep & \fChalTwo & \\
Hinglish B (synth.), F1 & \fHinBRep & \fHinBDep & \fHinBTwo & \pHinB \\
Hinglish B (synth.), FPR & \fprHinBRep & \fprHinBDep & \fprHinBTwo & \\
\textbf{Real Hindi, F1} & \fRHRep & \fRHDep & \fRHTwo & \pRH \\
\textbf{Real Hindi, FPR} & \fprRHRep & \fprRHDep & \fprRHTwo & \\\bottomrule
\end{tabular}
\end{table}

\paragraph{English.} On the real test set, v2 reaches \accRealTwo\% accuracy and \fRealTwo\% F1 (95\% CI
\ciLoRealTwo--\ciHiRealTwo), with ROC-AUC \aucRealTwo{} (Table~\ref{tab:main}). The difference from v1 deployed is not
significant ($p=\pReal$). The mismatch in v1 cost 0.6 F1 points and 0.8 points of false alarms against the reported
figure. On the challenge set, the appended summary happened to help v1, which shows that such a feature can look
useful while being inconsistent with training.

\begin{table}[t]\centering\small
\caption{Recall on scams (\%) on the attacked real test set. FPR on clean text: v1 \fprAdvCleanOne\%, v2 \fprAdvCleanTwo\%.}
\label{tab:adv}
\begin{tabular}{@{}lccc@{}}\toprule
Attack & msgs changed & v1 dep. & v2 \\\midrule
None & -- & \recCleanOne & \recCleanTwo \\
Homoglyph & \chgHomo & \recHomoOne & \recHomoTwo \\
Zero-width & \chgZero & \recZeroOne & \recZeroTwo \\
Spaced letters & \chgSpaced & \recSpacedOne & \recSpacedTwo \\
Leetspeak & \chgLeet & \recLeetOne & \recLeetTwo \\
Defanged URL & \chgDefang & \recDefangOne & \recDefangTwo \\
Combined & \chgCombo & \recComboOne & \recComboTwo \\\bottomrule
\end{tabular}
\end{table}

\paragraph{Evasion.} Each word-level trick costs v1 2.0--2.4 points of recall, defanged links cost 0.4, and the
combined attack costs 8.0.
v2 is essentially unaffected (Table~\ref{tab:adv}). Under attack, v1's false alarms fall slightly, because the
attacks also hide genuine messages' trigger words. The threat is real: the clean test set already contains ``Y0N0''
for YONO and ``H.D.F.C''. Character-level attacks are a known weakness of NLP models~\cite{boucher2022}, and here a
few lines of normalisation remove most of it.

\paragraph{Synthetic versus real Hindi.} v2 is perfect on synthetic set B and significantly better than v1
($p$ \pHinB). On real Hindi messages, every system is weak: v2 reaches \fRHTwo\% F1 (CI \ciLoRHTwo--\ciHiRHTwo) and
\fprRHTwo\% false alarms, against \fRHDep\% and \fprRHDep\% for v1 deployed. Group-level accuracy is \grpRHTwo\%. The
gain is not significant ($p=\pRH$). We see two causes:
\begin{itemize}\itemsep0pt
\item \emph{Missed scams.} Real lures are absent from the templates. Examples are government-scheme loans (``PM loan
  yojana, 1\% interest, call \dots''), absurd sales (``camera for Rs~199'') and multi-level-marketing pitches.
\item \emph{False alarms.} Genuine Devanagari notices are over-flagged, such as Covid test results, ration
  distribution messages and operator recharge reminders. Most Devanagari text in training is scam text, so the script
  itself has become evidence.
\end{itemize}
Templates written by the same people who wrote the rules measure coverage of those phrasings, not accuracy. We did
not tune anything on the real set.

\paragraph{Temporal split.} We trained without any campaign first seen in 2021 or later. The model then caught
\recFuture\% of \nFutureMsgs{} scam messages from \nFutureCamps{} campaigns first seen in 2022--23. Random hold-outs
of the same size scored \recRandMin--\recRandMax\% (mean \recRandMean\%). No drift is visible, but undated
campaigns remain in training, so this test is weak.

\paragraph{Calibration and the review band.} ECE is \eceClf{} for the classifier and \eceFinal{} for the final score.
Final scores in $[0.3, 0.7)$ cover \bandShare\% of real test messages and \bandErrIn{} of the \bandErrTotal{} errors.
Outside the band, accuracy is \bandAccOut\%. v2 marks such results ``not sure: check before you act'' rather than
forcing a label.

\begin{table}[t]\centering\small
\caption{Fusion ablation for v2 (F1 / FPR, \%).}
\label{tab:abl}
\begin{tabular}{@{}lccc@{}}\toprule
Variant & Real test & Challenge & Hinglish B \\\midrule
Classifier only & \ablRealClfF/\ablRealClfFpr & \ablChalClfF/\ablChalClfFpr & \ablHinBClfF/\ablHinBClfFpr \\
Rules only & \ablRealRulesF/\ablRealRulesFpr & \ablChalRulesF/\ablChalRulesFpr & \ablHinBRulesF/\ablHinBRulesFpr \\
+ intents & \ablRealIntentsF/\ablRealIntentsFpr & \ablChalIntentsF/\ablChalIntentsFpr & \ablHinBIntentsF/\ablHinBIntentsFpr \\
+ protective & \ablRealProtF/\ablRealProtFpr & \ablChalProtF/\ablChalProtFpr & \ablHinBProtF/\ablHinBProtFpr \\
Full (+ cap) & \ablRealFullF/\ablRealFullFpr & \ablChalFullF/\ablChalFullFpr & \ablHinBFullF/\ablHinBFullFpr \\\bottomrule
\end{tabular}
\end{table}

\paragraph{Fusion.} On the in-distribution real test set, the classifier alone is best (Table~\ref{tab:abl}). The
rules add false alarms, but they add 5.8 F1 points on the challenge set and 8.7 on Hinglish set B, at zero false
alarms on both. We keep the rules for that robustness and because they produce the explanations.

\paragraph{Simple rules.} ``Contains a link'' scores \linkRealF\% F1 with \linkRealFpr\% false alarms on the real
test set. It misses all \nNoLink{} scams without a link; v2 catches \recNoLink\% of them. The same rule scores
\linkCrossF\% on the India SMS set used in the v1 cross-domain study~\cite{scamguard1}. That benchmark's spam is
mostly link-bearing, so its scores should be read with this in mind.

\paragraph{Cost.} The full path takes \msPerMsg{} ms per message on two CPU cores.

\section{Limitations}
\begin{itemize}\itemsep0pt
\item The real Hindi set is small (\nRH{} messages) and was labelled by a single annotator.
\item The temporal test is weakened by undated campaigns.
\item Almost all Hindi and Hinglish training data is synthetic. A one-command ingest path splits newly collected real
  messages by group and keeps the frozen test set clean.
\item Multilingual transformers (MuRIL~\cite{muril}, XLM-R~\cite{xlmr}) on identical splits, including the real Hindi
  set, are prepared but not yet run. We report no transformer numbers.
\item No user study has been run.
\end{itemize}

\section{Ethics and reproducibility}
The no-login check and the Telegram channel never store messages. Phone numbers, account numbers, emails and UPI IDs
in the data are replaced by random stand-ins. From \texttt{backend/}, \texttt{python -m
ml\_training.research.v2\_experiments} regenerates every number in this paper, and this paper's numbers are
generated from its output.

\begin{thebibliography}{9}\small
\bibitem{scamguard1} P.~Thakur, B.~Goyal, S.~Chaudhary, S.~Subodh, and K.~N, ``ScamGuard: AI-based online scam and fraud detection system,'' under review, 2026.
\bibitem{imc25} ``Smishing-Dataset-IMC25,'' GitHub repository, \url{https://github.com/reportsmishing/Smishing-Dataset-IMC25}, CC BY 4.0.
\bibitem{sculley2015} D.~Sculley et al., ``Hidden technical debt in machine learning systems,'' in \emph{NeurIPS}, 2015.
\bibitem{mcnemar1947} Q.~McNemar, ``Note on the sampling error of the difference between correlated proportions or percentages,'' \emph{Psychometrika}, vol.~12, pp.~153--157, 1947.
\bibitem{guo2017} C.~Guo, G.~Pleiss, Y.~Sun, and K.~Q.~Weinberger, ``On calibration of modern neural networks,'' in \emph{ICML}, 2017.
\bibitem{boucher2022} N.~Boucher, I.~Shumailov, R.~Anderson, and N.~Papernot, ``Bad characters: Imperceptible NLP attacks,'' in \emph{IEEE S\&P}, 2022.
\bibitem{muril} S.~Khanuja et al., ``MuRIL: Multilingual representations for Indian languages,'' arXiv:2103.10730, 2021.
\bibitem{xlmr} A.~Conneau et al., ``Unsupervised cross-lingual representation learning at scale,'' in \emph{ACL}, 2020.
\end{thebibliography}
\end{document}
